\section{Integrating rational functions}

A rational function is a quotient of polynomials. Integrals of rational functions can range from simple to quite technical. In this course, the goal is to learn the basic reading habits: simplify first, look for a derivative of the denominator, and recognize standard logarithmic and arctangent patterns.

\subsection*{When the numerator is the derivative of the denominator}

If the integrand has the form
\[
\frac{g'(x)}{g(x)},
\]
then the integral is logarithmic:
\[
\int \frac{g'(x)}{g(x)}\,dx=\ln|g(x)|+C.
\]

\begin{examplebox}
Compute
\[
\int \frac{2x-1}{x^2-x+1}\,dx.
\]
The denominator is
\[
g(x)=x^2-x+1.
\]
Its derivative is
\[
g'(x)=2x-1.
\]
Therefore
\[
\int \frac{2x-1}{x^2-x+1}\,dx=\ln|x^2-x+1|+C.
\]
\end{examplebox}

The important step is recognizing the structure. The numerator is not random; it is the derivative of the denominator.

\subsection*{Adjusting the numerator}

Sometimes the numerator is close to the derivative of the denominator but not exactly equal. We can split it into a derivative part and a remainder.

\begin{examplebox}
Compute
\[
\int \frac{x+1}{x^2-x+1}\,dx.
\]
The derivative of the denominator is
\[
2x-1.
\]
Write the numerator \(x+1\) in terms of \(2x-1\):
\[
x+1=\frac12(2x-1)+\frac32.
\]
Thus
\[
\int \frac{x+1}{x^2-x+1}\,dx
=\frac12\int \frac{2x-1}{x^2-x+1}\,dx
+\frac32\int \frac{dx}{x^2-x+1}.
\]
The first integral is logarithmic:
\[
\frac12\ln|x^2-x+1|.
\]
For the second integral, complete the square:
\[
x^2-x+1=\left(x-\frac12\right)^2+\frac34.
\]
Therefore the remaining part leads to an arctangent form.
\end{examplebox}

This example shows the general idea. A rational integral may contain both logarithmic and arctangent pieces.

\subsection*{Completing the square}

The standard arctangent pattern is
\[
\int \frac{dx}{x^2+a^2}=\frac1a\arctan\frac{x}{a}+C,
\qquad a>0.
\]
When a quadratic denominator has no real roots, completing the square may reveal this pattern.

\begin{examplebox}
Compute
\[
\int \frac{dx}{x^2+4x+8}.
\]
Complete the square:
\[
x^2+4x+8=(x+2)^2+4.
\]
Thus
\[
\int \frac{dx}{x^2+4x+8}
=\int \frac{dx}{(x+2)^2+2^2}.
\]
Using the arctangent pattern,
\[
\int \frac{dx}{(x+2)^2+2^2}
=\frac12\arctan\frac{x+2}{2}+C.
\]
\end{examplebox}

\subsection*{Polynomial division}

If the degree of the numerator is greater than or equal to the degree of the denominator, divide first. Polynomial division separates the rational function into a polynomial part and a proper rational part.

\begin{examplebox}
Consider
\[
\int \frac{x^2+1}{x-1}\,dx.
\]
Divide:
\[
\frac{x^2+1}{x-1}=x+1+\frac2{x-1}.
\]
Therefore
\[
\int \frac{x^2+1}{x-1}\,dx
=\int (x+1)\,dx+2\int \frac{dx}{x-1}.
\]
Hence
\[
\int \frac{x^2+1}{x-1}\,dx
=\frac{x^2}{2}+x+2\ln|x-1|+C.
\]
\end{examplebox}

Simplifying before integrating is often the main step.

\begin{summarybox}
When integrating a rational function, ask:
\begin{itemize}
\item Is the fraction proper, or should I divide polynomials first?
\item Is the numerator the derivative of the denominator?
\item Can the numerator be split into a derivative part plus a remainder?
\item Does completing the square reveal an arctangent form?
\item Are logarithms expected from factors in the denominator?
\end{itemize}
\end{summarybox}

Rational integration is a large topic. The first skill is to read the denominator and numerator as related structures, not as isolated polynomials.